\subsection{EXAMPLE: POLYNOMIALS WITH INTEGER VALUES}

Let V be a finite dimensional \(\mathbf{Q}\)-vector space, \(V^{*}\) its dual, \(\mathcal{V}\) a lattice in V, \(\mathcal{V}^{*}\) the dual \(\mathbf{Z}\)-module of \(\mathcal{V}\), which can be identified canonically with a lattice in \(V^{*}\) , \(\mathbf{S}(V)\) the symmetric algebra of V, and

\[
\lambda : \mathbf {S} (\mathrm{V}) \to \mathrm{A} (\mathrm{V} ^ {*})
\]

the canonical bijection from \(\mathbf{S}(\mathrm{V})\) to the algebra of polynomial functions on \(\mathrm{V}^*\) (Algebra, Chap. IV, §5, no. 11, Remark 1). If we identify \(\mathrm{A}(\mathrm{V}^* \times \mathrm{V}^*)\) with \(\mathrm{A}(\mathrm{V}^*) \otimes_{\mathbf{Q}} \mathrm{A}(\mathrm{V}^*)\) , then \(\lambda\) transforms the coproduct of \(\mathbf{S}(\mathrm{V})\) into the map \(\mathrm{A}(\mathrm{V}^*) \to \mathrm{A}(\mathrm{V}^* \times \mathrm{V}^*)\) which associates to the polynomial function \(\varphi\) on \(\mathrm{V}^*\) the polynomial function

\[
(x, y) \mapsto \varphi (x + y)
\]

on \(\mathrm{V}^* \times \mathrm{V}^*\) (Algebra, Chap. IV, §5, no. 11, Remark 2).

Denote by \(\binom{\mathcal{V}}{\mathbf{Z}}\) the subset of \(\mathbf{S}(\mathrm{V})\) consisting of the elements which correspond to polynomial maps from \(\mathrm{V}^*\) to \(\mathbf{Q}\) that take integer values on \(\mathcal{V}^*\) .

PROPOSITION 2. (i) \(\begin{pmatrix} \mathcal{V} \\ \mathbf{Z} \end{pmatrix}\) is a biorder in the bigebra \(\mathbf{S}(\mathrm{V})\) , and \(\begin{pmatrix} \mathcal{V} \\ \mathbf{Z} \end{pmatrix} \cap \mathrm{V} = \mathcal{V}\) .

\begin{enumerate}
\def\labelenumi{(\roman{enumi})}
\setcounter{enumi}{1}
\item
  The \(\mathbf{Z}\) -algebra \(\left( \begin{array}{c}\mathcal{V}\\ \mathbf{Z} \end{array} \right)\) is generated by the \(\left( \begin{array}{c}h\\ n \end{array} \right)\) for \(h\in \mathcal{V},n\in \mathbf{N}\) .
\item
  If \((h_1, \ldots, h_r)\) is a basis of \(\mathcal{V}\) , the elements
\end{enumerate}

\[
\binom{h}{\mathbf {n}} = \binom{h _ {1}}{n _ {1}} \dots \binom{h _ {r}}{n _ {r}},
\]

where \(\mathbf{n} = (n_1, \ldots, n_r)\) belongs to \(\mathbf{N}^r\) , form a basis of the \(\mathbf{Z}\) -module \(\binom{\mathcal{V}}{\mathbf{Z}}\) .

For \(m \in \mathbf{N}\) , put \(\mathbf{S}_{m}(\mathrm{V}) = \sum_{i \leq m} \mathbf{S}^{i}(\mathrm{V}), \mathbf{S}_{m}(\mathcal{V}) = \sum_{i \leq m} \mathbf{S}^{i}(\mathcal{V})\) . By Algebra, Chap. IV, §5, no. 9, Prop. 15 and Remark,

\[
\mathbf {S} _ {m} (\mathcal {V}) \subset \mathbf {S} _ {m} (\mathrm{V}) \cap \binom{\mathcal {V}}{\mathbf {Z}} \subset \frac {1}{m !} \mathbf {S} _ {m} (\mathcal {V})
\]

so \(\binom{\mathcal{V}}{\mathbf{Z}}\cap\mathrm{V}=\mathcal{V}\) . Since \(\mathbf{S}_{m}(\mathcal{V})\) is a lattice in \(\mathbf{S}_{m}(\mathrm{V})\) , \(\mathbf{S}_{m}(\mathrm{V})\cap\binom{\mathcal{V}}{\mathbf{Z}}\) is also a lattice in \(\mathbf{S}_{m}(\mathrm{V})\) . On the other hand, \(\mathbf{S}_{m}(\mathrm{V})\cap\binom{\mathcal{V}}{\mathbf{Z}}\) is a direct factor of \(\mathbf{S}_{m+1}(\mathrm{V})\cap\binom{\mathcal{V}}{\mathbf{Z}}\) (since the quotient is torsion-free), hence it admits a complement which is a free \(\mathbf{Z}\)-module. It follows that \(\binom{\mathcal{V}}{\mathbf{Z}}\) is a free \(\mathbf{Z}\)-module. It is clear that this is a unital order in the algebra \(\mathbf{S}(\mathrm{V})\) . Let \((u_{n})_{n\in\mathbf{N}}\) be a basis of the \(\mathbf{Z}\)-module \(\binom{\mathcal{V}}{\mathbf{Z}}\) . This is also a basis of the \(\mathbf{Q}\)-module \(\mathbf{S}(\mathrm{V})\) and, for all

\[
\varphi \in \mathbf {S} (\mathrm{V} \times \mathrm{V}) = \mathbf {S} (\mathrm{V}) \otimes_ {\mathbf {Q}} \mathbf {S} (\mathrm{V}),
\]

there exists a unique sequence \((v_{n})\) of elements of S(V) such that \(\varphi=\sum u_{n}\otimes v_{n}\) . As above, identify S(V) with A(V*) and S(V) \(\otimes\) S(V) with A(V* × V*). If \(\varphi\in\binom{\mathcal{V}\times\mathcal{V}}{\mathbf{Z}}\) , the polynomial function \(x\mapsto\varphi(x,y)\) belongs to \(\binom{\mathcal{V}}{\mathbf{Z}}\) for all \(y\in\mathcal{V}^{*}\) . It follows that \(v_{n}(y)\in\mathbf{Z}\) for all n and all \(y\in\mathcal{V}^{*}\) , in other words that \(v_{n}\in\binom{\mathcal{V}}{\mathbf{Z}}\) . This proves that the coproduct maps \(\binom{\mathcal{V}}{\mathbf{Z}}\) to \(\binom{\mathcal{V}}{\mathbf{Z}}\otimes_{\mathbf{Z}}\binom{\mathcal{V}}{\mathbf{Z}}\) . If \(h\in\mathcal{V}\) and \(n\in\mathbf{N}\) , then \(\binom{h}{n}\) maps \(u\in\mathcal{V}^{*}\) to the integer \(\binom{u(h)}{n}\) , so \(\binom{h}{n}\in\binom{\mathcal{V}}{\mathbf{Z}}\) . Assertion (iii) is now obtained by applying Th. 1 to the commutative Lie algebra V, and (ii) follows.

COROLLARY. Let \(\mathrm{X}\) be an indeterminate. The polynomials \(\binom{\mathrm{X}}{n}\) , where \(n \in \mathbf{N}\) , form a basis of the \(\mathbf{Z}\) -module consisting of the polynomials \(\mathrm{P} \in k[\mathrm{X}]\) such that \(\mathrm{P}(\mathbf{Z}) \subset \mathbf{Z}\) .

If \(\mathrm{P}(\mathbf{Z}) \subset \mathbf{Z}\) , the Lagrange interpolation formula (Algebra, Chap. IV, §2, no. 1, Prop. 6) shows that the coefficients of P belong to \(\mathbf{Q}\) ; thus, we can assume that \(k = \mathbf{Q}\) and apply Prop. 2 with \(V = \mathbf{Q}\) , \(\mathcal{V} = \mathbf{Z}\) .

